\section{What the Literature Review Must Establish}
This review examines four connected questions: what earlier EEG work says about dementia classification; how recording duration and validation affect conclusions; what pretrained EEG encoders and hybrid decoders contribute; and why accuracy needs to be read alongside reliability and uncertainty. The comparison table includes 30 papers selected for those purposes. They span established EEG signal studies, work directly using the public Alzheimer's cohort, foundation-model papers and methods papers needed to justify subject-disjoint evaluation. The latest papers available by November 2026 are included. Results from different papers are \emph{not} treated as directly comparable benchmarks because they often use different samples, preprocessing choices and validation rules.

\section{EEG Signals and Alzheimer's Classification}
Early studies associate Alzheimer's disease with changes in spectral activity and dynamical structure in EEG \cite{jeong2004,dauwels2010}. These observations motivate frequency-band features, including delta, theta, alpha, beta and gamma power. They also warn against treating a single measure as a disease-specific signal. Review work highlights substantial variation in how recordings are cleaned, features are computed and patient groups are separated \cite{cassani2018,tzimourta2021review}.

The public dataset used here contains EEG recordings from 36 AD participants, 23 with frontotemporal dementia and 29 cognitively normal controls \cite{miltiadous2023dataset}. Published analyses of this collection include classical feature combinations and convolution-transformer designs \cite{zheng2023,dice2023}. These studies motivate comparison with both spectral baselines and a more flexible representation model. However, their published accuracy values should not be read as head-to-head evidence against our results unless the participants, data partitioning and testing conditions match.

\section{Recording Duration and Participant-Level Validation}
EEG window duration has already been investigated specifically for Alzheimer's classification. Tzimourta and colleagues evaluated window length and regional differences in 2019 \cite{tzimourta2019}. Our temporal contribution is therefore not the first use of recording length as a variable. It instead asks how prediction correctness and within-recording consistency behave under a stated subject-level evaluation, and how stable but wrong decisions complicate the meaning of ``reliability.''

Short EEG segments from one participant are correlated. If some segments train a model and others test it, the apparent sample size becomes misleading for diagnosing previously unseen people. Nested or held-out model selection is important because choosing parameters on the same examples used for final assessment can bias results \cite{varma2006,varoquaux2018}. The 2026 AHEPA review found substantially different accuracy patterns when studies were grouped by validation rigor \cite{miltiadous2026benchmark}. This supports our choice to keep participants, not windows, as the independent evaluation unit.

\section{Pretrained EEG Encoders and Custom Decoders}
Deep convolutional EEG models can learn short-range features from signals without hand-selecting all features \cite{schirrmeister2017,lawhern2018}. Hybrid convolution-transformer models can combine local and longer-range information \cite{song2023,dice2023}. More recently, BENDR, BIOT, LaBraM and EEGPT have explored reusable representations learned from many biosignal recordings \cite{kostas2021,yang2023biot,jiang2024labram,wang2024eegpt}. CBraMod, published at ICLR 2025, models spatial and temporal EEG dependencies through a criss-cross architecture and provides the pretrained encoder used in this project \cite{wang2025cbramod}.

A pretrained backbone does not determine how best to classify one small downstream cohort. Our DSRF decoder combines high-level pretrained features with low-level temporal convolution features, then uses a learned gate, a residual path and class-attention pooling. The idea is related to broader local/global fusion methods, not established as a completely new category of neural network. The relevant test is whether this implementation remains competitive under comparisons where participants are separated and the final holdout is not repeatedly used for model selection.

The literature developed rapidly in 2026. LEAD studied an EEG foundation model specifically for Alzheimer's disease, while EEG-FM-Bench and EEG-FM-Compass stressed systematic benchmarking and adaptation choices \cite{wang2026lead,xiong2026,liu2026compass}. Their existence makes it especially important to frame EEGraft as a targeted decoder and reliability study rather than the first Alzheimer's EEG foundation model.

\section{Comparative Literature Table}
Table~\ref{tab:literature} compares 30 relevant publications. Its final column explains what each paper changes about our design, evidence standard or interpretation. A paper's reported result is not presented as a result of this thesis.
\begingroup\footnotesize
\setlength{\tabcolsep}{3.3pt}\renewcommand{\arraystretch}{1.16}
\begin{longtable}{@{}>{\RaggedRight\hyphenpenalty=3000\exhyphenpenalty=4500\arraybackslash}p{0.185\textwidth}>{\RaggedRight\hyphenpenalty=3000\exhyphenpenalty=4500\arraybackslash}p{0.305\textwidth}>{\RaggedRight\hyphenpenalty=3000\exhyphenpenalty=4500\arraybackslash}p{0.40\textwidth}@{}}
\caption{Thirty representative research papers and their relevance to EEGraft.}\label{tab:literature}\\
\toprule
\textbf{Paper} & \textbf{Main focus} & \textbf{Relevance and limitation for \mbox{EEGraft}}\\
\midrule\endfirsthead
\caption[]{Literature comparison (continued)}\\
\toprule\textbf{Paper} & \textbf{Main focus} & \textbf{Relevance and limitation for \mbox{EEGraft}}\\\midrule\endhead
\bottomrule\endfoot
Jeong (2004) \cite{jeong2004}& EEG dynamics in Alzheimer's disease & Motivates spectral and dynamical analysis; descriptive EEG effects are not a stand-alone diagnosis.\\
Dauwels et al. (2010) \cite{dauwels2010}& Review of slowing, complexity and synchrony & Supports multiple signal descriptors and careful artifact handling.\\
Cassani et al. (2018) \cite{cassani2018}& Systematic review of resting-state EEG for AD & Shows variation in algorithms and evidence quality; warns against direct cross-paper accuracy rankings.\\
Tzimourta et al. (2019) \cite{tzimourta2019}& EEG window length and regional analysis & Establishes prior work on duration; our claim concerns stability versus correctness, not first study of window length.\\
Tzimourta et al. (2021) \cite{tzimourta2021review}& Review of AD EEG learning methods & Grounds traditional baselines and study-design challenges.\\
Miltiadous et al. (2021) \cite{miltiadous2021}& AD/FTD EEG classification and validation choices & Directly relevant to our disease tasks and validation discipline.\\
Miltiadous et al. (2023) \cite{miltiadous2023dataset}& Public 88-subject AD/FTD/CN EEG dataset & Defines cohort provenance and available 19-channel resting EEG; not evidence of external generalization.\\
Miltiadous et al. (2023), DICE-Net \cite{dice2023}& Convolution-transformer architecture for dementia EEG & Closely related hybrid architecture; outcome values require protocol-aware comparison.\\
Zheng et al. (2023) \cite{zheng2023}& Joint spectral, complexity and synchrony features & Supports explicit classical feature baselines and subject-oriented evaluation.\\
Puri et al. (2024) \cite{puri2024}& Lightweight spatiotemporal CNN (LEADNet) & Shows why smaller networks remain worth comparing with pretrained models.\\
Alidadi et al. (2026) \cite{alidadi2026}& Review of past and current AD EEG methods & Highlights persistent generalization, bias and explanation challenges.\\
Miltiadous et al. (2026) \cite{miltiadous2026benchmark}& AHEPA study-validity scoping review & Directly reinforces subject-wise splits and warns about inflated epoch-wise accuracy.\\
Masry et al. (2026) \cite{masry2026}& Montage-agnostic federated EEG AD analysis & Emphasizes channel mismatch, privacy and the difference between simulated and true site transfer.\\
Akbar et al. (2026) \cite{akbar2026}& Feature engineering, deep SVM and explanations & Another hybrid dementia method; attribution is model-related, not automatically causal evidence.\\
Schirrmeister et al. (2017) \cite{schirrmeister2017}& CNN feature learning from EEG & Basis for learning temporal filters directly from signals.\\
Lawhern et al. (2018) \cite{lawhern2018}& Compact EEGNet architecture & Motivates a small neural baseline rather than only comparing complex models.\\
Song et al. (2023) \cite{song2023}& EEG Conformer local/global representation & Supports combining convolutional and attention mechanisms; not Alzheimer's-specific proof.\\
Kostas et al. (2021) \cite{kostas2021}& Self-supervised BENDR EEG encoding & Precedent for transferring pretrained EEG representations to downstream tasks.\\
Yang et al. (2023) \cite{yang2023biot}& BIOT, biosignal transfer with varying formats & Explains potential of flexible pretraining, but transfer may remain task-dependent.\\
Jiang et al. (2024) \cite{jiang2024labram}& LaBraM pretrained EEG representation & Alternative to task-specific training; does not establish usefulness on our cohort.\\
Wang et al. (2024) \cite{wang2024eegpt}& EEGPT representation alignment & Relevant to robustness of pretrained representations and adapter choice.\\
Wang et al. (2025) \cite{wang2025cbramod}& CBraMod criss-cross foundation model & The actual pretrained backbone in EEGraft; it is not our invention.\\
Wang et al. (2026), LEAD \cite{wang2026lead}& Alzheimer's-specific EEG foundation model & Makes the foundation-model comparison more current; we did not benchmark LEAD on the locked cohort.\\
Xiong et al. (2026) \cite{xiong2026}& EEG-FM-Bench cross-model evaluation & Supports standardized protocols and cautious representation comparisons.\\
Liu et al. (2026) \cite{liu2026compass}& EEG-FM-Compass review and benchmark & Frames adaptation, model scale and transfer as open evaluation questions.\\
Delorme and Makeig (2004) \cite{delorme2004}& EEGLAB signal processing and ICA & Connects public EEG preprocessing provenance with accepted signal-processing tools.\\
Pernet et al. (2019) \cite{pernet2019}& EEG-BIDS data organization & Supports auditable participant identities, channel metadata and recording conditions.\\
Chen and Guestrin (2016) \cite{chen2016xgboost}& XGBoost tree boosting & Defines an important classical competitor that later led the locked test.\\
Varma and Simon (2006) \cite{varma2006}& Model-selection bias in cross-validation & Explains why development tuning and final testing must be separated.\\
Varoquaux (2018) \cite{varoquaux2018}& Uncertainty with small neuroimaging samples & Explains why a 13-person holdout cannot establish a small performance difference reliably.\\
\end{longtable}
\enlargethispage{3\baselineskip}\endgroup

\section{Synthesis and Research Gap}
Three conclusions emerge. First, EEG disease classification has strong historical grounding, but its results depend on sampling, preprocessing and evaluation design. Second, changing window length and model architecture is not new on its own; a defensible contribution requires a precise question and trustworthy comparisons. Third, a more sophisticated representation does not eliminate the need for baselines, locked testing, calibration checks and failure analysis \cite{miltiadous2026benchmark,xiong2026,guo2017calibration}.

Our research gap is therefore narrower and more useful than a claim to invent the first EEG Alzheimer's detector. We examine how often apparent accuracy conflicts with stability and correctness, and then assess a custom frozen-backbone decoder through subject-level comparisons and stress tests. This directly motivates the two-track design described in Chapter 4.\par
